% ID: SEL-D0055-OB-E03
% ESTADO: cotejada. Segunda pasada visual y borrador compilado revisados.
% Fuente: Selectividad/sel_2011_matematicas_aplicadas/122A-EXAMEN 6.pdf
% SHA-256: aebcbfddd4d5db5f5a10b4662e9e12b2ee362bbf574b751bccd6d78594766ad3
% Página: 2; opción B; ejercicio original 3.

Pedro vive en una ciudad donde el 40\% de los días del año hay riesgo de lluvia y el resto no lo hay. Cuando hay riesgo de lluvia, Pedro coge el paraguas un 98\% de las veces y cuando no lo hay, un 5\% de las veces. Si se selecciona un día del año al azar,
\begin{enumerate}[label=\alph*)]
\item \textbf{(1.25 puntos)} \textquestiondown{}cuál es la probabilidad de que Pedro no haya cogido el paraguas ese día?
\item \textbf{(1.25 puntos)} \textquestiondown{}cuál es la probabilidad de que exista riesgo de lluvia, si sabemos que ese día Pedro ha cogido el paraguas?
\end{enumerate}
